
\section{Discussion}
\label{section:Discussion}

\paragraph{Verification in the Clifford+MSI model.}
The central contribution of this paper is to provide a framework for verification protocols that simultaneously achieves composable security, exponentially small soundness error against malicious behavior, and robustness to circuit-level noise, while being tailored to the Clifford+MSI model.
Compared to the recent result of \cite{BN25noise}, which focuses on noise robustness for verification in the circuit-model, our work additionally provides a composable and modular construction for verification protocols.
This closes the conceptual and practical gaps that had been widening between circuit- and MBQC-model-based verification strategies.

\paragraph{Constructions for composable delegation.}
Along the way, we introduced a delegation protocol that is sufficient to hide whether a computation or a test is being delegated, with full composable security inherited from its sub-components. Working in the Clifford+MSI model, we derived a blind state-injection protocol that hides the potential magic injected by the Client at each layer, together with a protocol that delegates a Clifford circuit and measurements on a blinded quantum state. Since these primitives are composably secure by construction, they can be reused in other protocols and optimized independently, both theoretically and in practical implementations.

\paragraph{A framework for verification in the circuit-model.}
By generalizing the trap design, we have shown that there exists a unified stabilizer-based formalism for circuit-model verification. It allows trap designs on subsets of qubits (rather than single qubits) and supports the combination of test rounds by merging compatible traps. This modularity exposes a whole family of verification protocols—including \cite{B18how,BN25noise}—all enjoying the same security guarantees. This has the potential to enable Clients to derive trap-based protocols optimized for noise-robustness (see paragraph below) and to tailor verification schemes to hardware constraints, in direct analogy with dummyless verification in MBQC \cite{KKLM23asymmetric}.

\paragraph{Noise-robustness and trap engineering.}
We highlight that the verification framework presented here achieves noise-robustness for circuit-level noise, adopting the same noise model as established in \cite{LMKO21verifying, BN25noise}. In our construction, the maximum tolerated noise rate is bounded by $\alpha \cdot \rate$, where $\alpha$ is a constant related to the $\BQP$ error and $\rate$ represents the \textit{trap detection rate} (i.e., the probability that a harmful Pauli-error is detected by at least one of the traps). In the current case, which is the simplest, the protocol chooses uniformly between $k$ types of test rounds, and thus this rate is $\rate = 1/k$. For the specific construction of Protocol \ref{protocol:verification}, we have $\rate = 1/(1+t)$, where $t$ is the number of state injection layers, so the maximum tolerated noise rate is bounded by $\alpha/(1+t)$. On the other hand, the protocols in \cite{BN25noise} and \cite{LMKO21verifying} use two types of test runs only, so they achieve $k=2$, meaning $\rate=1/2$, and tolerate noise up to $\alpha/2$.

This contrast suggests that noise-robustness can be significantly improved by engineering traps to increase $\rate$. While a direct way to improve the rate is to reduce the number of test types $k$---a task made possible by the flexibility of our framework as teased in Section \ref{subsection:framework generalized}---our modular approach enables even more sophisticated optimizations.
Specifically, the detection rate $\rate$ does not need to be restricted to the naive $1/k$ lower bound derived from a uniform choice of test rounds. As showcased in \cite{KKLM22unifying} for the MBQC model, within the stabilizer formalism $\rate$ can be determined through a more refined analysis of the trap-based construction's detection capabilities. By applying similar analytical tools to our circuit-model framework, one could potentially improve noise-robust verification to meet the $\rate=1/2$ detection rate of \cite{BN25noise} without relying on the compilation trick of \cite{B18how} that introduces significant ancilla overhead. As a crucial consequence, this would achieve the same level of noise-tolerance more efficiently.

\paragraph{Magic-blindness is sufficient for verification. Is it necessary?}
We established \emph{magic-blindness} as a sufficient notion of blindness for
 verification in the circuit-model. This raises a more fundamental question: \emph{must a verification protocol necessarily hide some quantum computational resource from the prover?} In our construction, the hidden resource is \emph{magic} (non-stabilizerness), because it is what enables quantum advantage in Clifford+MSI architectures, but there is no a priori reason why this must be the only possibility. This reframes verification as a \emph{resource-hiding task}: is there a different, or minimal, quantum resource whose blindness is required and sufficient to make malicious deviations detectable? If such a minimal resource exists, it would imply fundamental lower bounds on verification overhead. More generally, this question opens the door to a \emph{resource theory of verification}.
